\chapter{Méthodes bayésiennes}\label{ch:qm:bayesian-methods}

Une plateforme multigérants examine ses cinquante gérants de portefeuille. Le meilleur \omterm{def:qm:estimation:sharpe}{ratio de Sharpe} sur trois ans est 2.4, et le comité d'allocation
propose de doubler le capital de ce gérant. Trois années permettent d'estimer un \omterm{def:qm:estimation:sharpe}{ratio de Sharpe} avec une \omterm{def:qm:estimation:estimator}{erreur type} d'environ 0.58, et le meilleur parmi cinquante
historiques bruités est, presque par construction, un bon gérant qui a aussi eu de la chance. Rétréci vers ce que les gérants de la plateforme obtiennent habituellement,
ce même historique vaut un \omterm{def:qm:estimation:sharpe}{ratio de Sharpe} de 1.14~; les trois années suivantes du gérant, dans la simulation du chapitre, donnent 1.69, et sur deux
mille plateformes simulées, le meilleur historique a une moyenne de 2.09 tandis que les trois années suivantes du même gérant ont une moyenne de 1.02. Ce chapitre porte sur
le raisonnement à partir de connaissances préalables~: \omterm{def:qm:bayesian-methods:prior}{lois a priori}, \omterm{def:qm:bayesian-methods:prior}{lois a posteriori} et familles conjuguées qui les rendent faciles~; le rétrécissement et l'\omterm{def:qm:estimation:estimator}{estimateur} qui
montre que la moyenne empirique peut être battue~; les \omterm{def:qm:bayesian-methods:hier}{modèles hiérarchiques} dont la \omterm{def:qm:bayesian-methods:prior}{loi a priori} est apprise à partir de la population~; et les échantillonneurs à \omterm{def:qm:markov-chains-and-queues:ctmc}{chaînes de Markov} qui
calculent les \omterm{def:qm:bayesian-methods:prior}{lois a posteriori} lorsqu'aucune formule ne le permet.

\section{Lois a priori, lois a posteriori et conjugaison}

\begin{definition}[Loi a priori, loi a posteriori, loi a priori conjuguée]\label{def:qm:bayesian-methods:prior}
Dans un modèle bayésien, le paramètre $\theta$ est aléatoire. Sa \emph{loi a priori}\index{loi a priori} $p(\theta)$ énonce ce que l'on croit
avant les données~; la \emph{loi a posteriori}\index{loi a posteriori} $p(\theta \mid x) \propto p(x \mid \theta)\,
p(\theta)$ est la mise à jour bayésienne de celle-ci par la vraisemblance des données $x$ (Bayes, 1763). Une famille de lois a priori est une \emph{loi a priori
conjuguée}\index{loi a priori conjuguée} pour une vraisemblance si la loi a posteriori reste dans la famille.
\end{definition}

\begin{definition}[Loi prédictive a posteriori, intervalle de crédibilité]\label{def:qm:bayesian-methods:predictive}
La \emph{loi prédictive a posteriori}\index{loi prédictive a posteriori} de nouvelles données $\tilde x$ est $p(\tilde x \mid x) = \int
p(\tilde x \mid \theta)\,p(\theta \mid x)\,d\theta$. Un \emph{intervalle de crédibilité}\index{intervalle de crédibilité} $(1 - \alpha)$ est un intervalle
auquel la \omterm{def:qm:bayesian-methods:prior}{loi a posteriori} attribue une probabilité $1 - \alpha$, en général entre ses quantiles d'ordre $\alpha/2$ et $1 - \alpha/2$.
\end{definition}

Contrairement à un intervalle de confiance, un \omterm{def:qm:bayesian-methods:predictive}{intervalle de crédibilité} est un énoncé de probabilité sur le paramètre, conditionnellement aux données, et il ne vaut pas plus que
la \omterm{def:qm:bayesian-methods:prior}{loi a priori}. Trois couples conjugués couvrent l'essentiel des usages dans les desks. Un taux de réussite $p$ muni d'une \omterm{def:qm:bayesian-methods:prior}{loi a priori} $\mathrm{Beta}(a, b)$, après $k$ succès en $n$ essais,
admet pour \omterm{def:qm:bayesian-methods:prior}{loi a posteriori} $\mathrm{Beta}(a + k, b + n - k)$~: un signal juste dans 110 des 200 transactions, avec une \omterm{def:qm:bayesian-methods:prior}{loi a priori} $\mathrm{Beta}(50, 50)$ qui traduit
l'idée que les signaux s'éloignent rarement d'une pièce équilibrée, a une moyenne a posteriori de 0.533, un \omterm{def:qm:bayesian-methods:predictive}{intervalle de crédibilité} à 95~\% $(0.477, 0.589)$ et une probabilité a posteriori de
0.876 que son taux de réussite dépasse la moitié. Une moyenne gaussienne avec une \omterm{def:qm:bayesian-methods:prior}{loi a priori} gaussienne donne une \omterm{def:qm:bayesian-methods:prior}{loi a posteriori} gaussienne (voir plus bas). Un échantillon gaussien
de moyenne et de variance inconnues admet pour loi conjuguée la famille normale-gamma inverse.

\begin{proposition}[La loi a posteriori normale-normale est une moyenne pondérée par les précisions]\label{prop:qm:bayesian-methods:normal}
Si $\theta \sim \mathcal N(m, \tau^2)$ et $x \mid \theta \sim \mathcal N(\theta, \sigma^2)$, alors $\theta \mid x \sim \mathcal N(\hat\theta,
v)$ avec $1/v = 1/\tau^2 + 1/\sigma^2$ et
\[
\hat\theta = \frac{m/\tau^2 + x/\sigma^2}{1/\tau^2 + 1/\sigma^2} = m + (1 - B)(x - m), \qquad B = \frac{\sigma^2}{\sigma^2 + \tau^2}.
\]
\end{proposition}

\begin{proof}
Le logarithme de la \omterm{def:qm:bayesian-methods:prior}{loi a posteriori} vaut $-\frac{(\theta - m)^2}{2\tau^2} - \frac{(x - \theta)^2}{2\sigma^2}$ à une constante près, soit un trinôme en $\theta$ de
coefficient $-\frac12(1/\tau^2 + 1/\sigma^2)$ sur $\theta^2$ et $m/\tau^2 + x/\sigma^2$ sur $\theta$~; en complétant le carré on obtient la moyenne
et la variance.
\end{proof}

Le facteur de rétrécissement $B$ est la part de l'écart observé $x - m$ attribuée au bruit. Il est grand lorsque les données sont bruitées
($\sigma$ grand) ou la population homogène ($\tau$ petit).

\section{Rétrécissement}

\begin{definition}[Estimateur à rétrécissement, estimateur de James--Stein]\label{def:qm:bayesian-methods:js}
Un \emph{estimateur à rétrécissement}\index{estimateur à rétrécissement} ramène un ensemble d'estimations vers une cible commune. Pour $X \sim \mathcal N(\theta,
\sigma^2I_d)$, l'\emph{estimateur de James--Stein}\index{estimateur de James--Stein} vers zéro vaut $\bigl(1 - (d - 2)\sigma^2/\lvert X\rvert^2\bigr)X$~;
vers la moyenne générale $\bar X$, il vaut $\bar X + \bigl(1 - (d - 3)\sigma^2/\sum_i(X_i - \bar X)^2\bigr)(X - \bar X)$, et sa partie positive
tronque le facteur à zéro.
\end{definition}

\begin{theorem}[L'estimateur de James--Stein domine la moyenne empirique]\label{thm:qm:bayesian-methods:js}
Pour $d \ge 3$ et tout $\theta$, l'\omterm{def:qm:bayesian-methods:js}{estimateur de James--Stein} vers zéro a une erreur quadratique moyenne totale strictement inférieure à celle de $X$~:
$\E\lvert\delta_{\mathrm{JS}} - \theta\rvert^2 = d\sigma^2 - (d - 2)^2\sigma^4\,\E[1/\lvert X\rvert^2] < d\sigma^2$.
\end{theorem}

\begin{proof}
Prenons $\sigma = 1$ et écrivons $\delta = X - g(X)$ avec $g(x) = (d - 2)x/\lvert x\rvert^2$. Alors $\E\lvert\delta - \theta\rvert^2 = d -
2\E[(X - \theta)\cdot g(X)] + \E\lvert g(X)\rvert^2$. L'identité de Stein, obtenue par intégration par parties contre la densité gaussienne, donne
$\E[(X_i - \theta_i)g_i(X)] = \E[\partial_ig_i(X)]$, et $\sum_i\partial_ig_i(x) = (d - 2)^2/\lvert x\rvert^2$. Comme $\lvert g(x)\rvert^2 = (d
- 2)^2/\lvert x\rvert^2$ aussi, le risque vaut $d - (d - 2)^2\E[1/\lvert X\rvert^2]$, et l'espérance est finie pour $d \ge 3$.
\end{proof}

Le résultat (James et Stein, 1961) a bouleversé la profession~: en estimant trois moyennes sans rapport, les moyennes empiriques sont inadmissibles. Le gain
est grand lorsque les moyennes sont proches de la cible relativement au bruit, ce qui est la situation des \omterm{def:qm:estimation:sharpe}{ratios de Sharpe} de cinquante gérants mesurés
sur trois ans. Efron et Morris (1975) l'ont rendu célèbre avec les moyennes au bâton.

\section{Modèles hiérarchiques et bayésien empirique}

\begin{definition}[Modèle hiérarchique, bayésien empirique]\label{def:qm:bayesian-methods:hier}
Un \emph{modèle hiérarchique}\index{modèle hiérarchique} attribue à de nombreuses unités une \omterm{def:qm:bayesian-methods:prior}{loi a priori} commune dont les paramètres propres sont
inconnus~: $x_i \mid \theta_i \sim \mathcal N(\theta_i, \sigma_i^2)$, $\theta_i \sim \mathcal N(m, \tau^2)$. Le \emph{bayésien
empirique}\index{bayésien empirique} estime $(m, \tau^2)$ à partir des données, en maximisant la vraisemblance marginale $x_i \sim \mathcal N(m, \sigma_i^2 +
\tau^2)$ ou par moments, puis applique à chaque unité les formules a posteriori en y reportant les estimations.
\end{definition}

Lorsque $\sigma_i = \sigma$ pour toutes les unités, les estimations par moments valent $\hat m = \bar x$ et $\hat\tau^2 = \max(0, s_x^2 - \sigma^2)$~: la dispersion des
historiques au-delà de ce que le seul bruit produirait. L'\omterm{def:qm:estimation:estimator}{estimateur} qui en résulte est celui de James--Stein à un facteur $(k - 3)/(k - 1)$ près, ce qui explique
le théorème précédent~: le rétrécissement vers la moyenne générale est un \omterm{def:qm:estimation:estimator}{estimateur} bayésien dont la \omterm{def:qm:bayesian-methods:prior}{loi a priori} est estimée à partir des autres unités.

\begin{example}[Cinquante gérants]\label{ex:qm:bayesian-methods:platform}
Les cinquante gérants de la plateforme ont des \omterm{def:qm:estimation:sharpe}{ratios de Sharpe} annuels véritables tirés de $\mathcal N(0.5, 0.4^2)$, et chaque historique sur trois ans ajoute un bruit
d'écart type $\sigma = 1/\sqrt3 = 0.577$ (chapitre~11). Sur la plateforme simulée à graine fixée, dont le meilleur historique vaut 2.41 comme dans l'accroche, les historiques
ont une moyenne de 0.51 et un écart type de 0.71, donc $\hat\tau = \sqrt{0.71^2 - 1/3} = 0.41$ et $B = 0.333/(0.333 + 0.168) = 0.665$~: deux tiers de
chaque écart à la moyenne est du bruit. Le meilleur historique devient $0.51 + 0.335 \times (2.41 - 0.51) = 1.14$, avec un écart type
a posteriori de 0.33~; l'\omterm{def:qm:bayesian-methods:js}{estimateur de James--Stein} donne 1.20. Les deuxième et troisième historiques, 1.82 et 1.45, deviennent 0.95 et 0.83~; le classement ne change pas, seules les
distances changent. Les ratios véritables des trois premiers sont 1.07, 0.96 et 1.68, et leurs trois années suivantes donnent 1.69, 0.57 et 0.99
(\cref{fig:qm:bayesian-methods:shrink}).
\end{example}

\begin{omfigure}
\begin{tikzpicture}
\begin{axis}[width=9.5cm,height=6.2cm,xlabel={Ratio de Sharpe, trois premières années},ylabel={Ratio de Sharpe, trois années suivantes},
  tick label style={font=\small},label style={font=\small},xmin=-1,xmax=3,ymin=-1,ymax=3,grid=major,grid style={gray!25},
  legend columns=3,legend style={font=\footnotesize,at={(0.5,-0.25)},anchor=north,draw=none,fill=none,
  /tikz/every even column/.append style={column sep=8pt}},table/col sep=comma]
\addplot[only marks,mark=*,mark size=1.6pt,omDef] table[x=x1,y=x2]{figdata/methods/14-bayesian-methods/managers.csv};
\addplot[black,dashed,thick] table[x=x,y=diag]{figdata/methods/14-bayesian-methods/lines.csv};
\addplot[omThm,very thick] table[x=x,y=eb]{figdata/methods/14-bayesian-methods/lines.csv};
\legend{gérants,sans rétrécissement,bayésien empirique}
\end{axis}
\end{tikzpicture}

\omcaption{Les \omterm{def:qm:estimation:sharpe}{ratios de Sharpe} des cinquante gérants sur deux périodes consécutives de trois ans. Prendre le premier historique au pied de la lettre prédit la
diagonale en tirets~; la moyenne a posteriori du \omterm{def:qm:bayesian-methods:hier}{bayésien empirique} prédit la droite plus plate, de pente 0.335 passant par la moyenne de la plateforme 0.51, et la seconde
période se disperse autour d'elle. Données~: le tutoriel du chapitre, avec graine fixée.}\label{fig:qm:bayesian-methods:shrink}
\end{omfigure}

Par rapport aux ratios vrais, le \omterm{def:qm:bayesian-methods:hier}{bayésien empirique} réduit l'erreur quadratique moyenne des cinquante estimations de 0.335 à 0.149 sur cette plateforme.
En moyenne sur 2\,000 plateformes simulées, il la ramène de 0.336 à 0.121, proche du 0.109 de l'oracle qui connaît $m$ et $\tau$~; et
face aux historiques des trois années suivantes, ce qu'un allocateur peut vérifier, l'erreur passe de 0.672 à 0.457, soit une réduction de 32~\%
(voir la \cref{fig:qm:bayesian-methods:mse}).

\begin{omfigure}
\begin{tikzpicture}
\begin{axis}[width=10cm,height=5cm,ybar,bar width=12pt,
  xtick={0,1,2,3},xticklabels={{historique\\brut},{James--\\Stein},{bayésien\\empirique},{oracle\\bayésien}},
  xticklabel style={align=center},
  ylabel={erreur quadratique moyenne},tick label style={font=\small},label style={font=\small},
  xmin=-0.6,xmax=3.6,ymin=0,ymax=0.8,grid=major,grid style={gray!25},
  legend style={font=\footnotesize,at={(0.97,0.97)},anchor=north east,fill=white,draw=none},table/col sep=comma]
\addplot[fill=omDef!60,draw=omDef] table[x=k,y=truth]{figdata/methods/14-bayesian-methods/mse.csv};
\addplot[fill=omThm!45,draw=omThm] table[x=k,y=next]{figdata/methods/14-bayesian-methods/mse.csv};
\legend{face aux ratios vrais,face aux trois années suivantes}
\end{axis}
\end{tikzpicture}

\omcaption{Erreur quadratique moyenne de quatre estimations des \omterm{def:qm:estimation:sharpe}{ratios de Sharpe} de cinquante gérants, en moyenne sur 2\,000 plateformes simulées~: les
historiques bruts sur trois ans (0.336 face à la vérité, 0.672 face à la période suivante), James--Stein et \omterm{def:qm:bayesian-methods:hier}{bayésien empirique} (0.121 et 0.457), et la
moyenne a posteriori avec les vrais paramètres de la plateforme (0.109 et 0.445). Données~: le tutoriel du chapitre, avec graine fixée.}\label{fig:qm:bayesian-methods:mse}
\end{omfigure}

Le meilleur historique est là où le rétrécissement compte le plus. Sur les 2\,000 plateformes, le meilleur de cinquante historiques de trois ans vaut en moyenne 2.09, tandis que
le \omterm{def:qm:estimation:sharpe}{ratio de Sharpe} vrai de ce même gérant vaut en moyenne 1.01~: la sélection seule ajoute 1.08. L'estimation bayésienne empirique du gérant sélectionné
vaut en moyenne 1.01, et les trois années suivantes 1.02 (voir la \cref{fig:qm:bayesian-methods:curse}). Le problème du comité relève d'une malédiction du vainqueur, soit
l'effet de sélection de la recherche du chapitre~12, et la \omterm{def:qm:bayesian-methods:prior}{loi a priori} en est la correction.

\begin{omfigure}
\begin{tikzpicture}
\begin{axis}[width=10cm,height=5cm,xlabel={estimation moins le ratio de Sharpe vrai, meilleur des cinquante},ylabel={densité},
  tick label style={font=\small},label style={font=\small},xmin=-1.5,xmax=3,ymin=0,ymax=1.6,grid=major,grid style={gray!25},
  legend style={font=\footnotesize,at={(0.97,0.97)},anchor=north east,fill=white,draw=none},table/col sep=comma]
\addplot[omThm,thick,const plot mark mid,fill=omThm!15] table[x=x,y=raw]{figdata/methods/14-bayesian-methods/curse.csv};
\addplot[omDef,very thick,const plot mark mid] table[x=x,y=eb]{figdata/methods/14-bayesian-methods/curse.csv};
\addplot[black,dashed,thin,sharp plot] coordinates {(0,0) (0,1.6)};
\legend{historique brut,bayésien empirique}
\end{axis}
\end{tikzpicture}

\omcaption{L'erreur sur le \omterm{def:qm:estimation:sharpe}{ratio de Sharpe} estimé du meilleur des cinquante gérants, sur 2\,000 plateformes simulées~: la
valeur historique brute surestime en moyenne de 1.08 le gérant retenu~; l'estimation bayésienne empirique est centrée sur la vérité. Données~: le tutoriel du chapitre,
avec graine fixée.}\label{fig:qm:bayesian-methods:curse}
\end{omfigure}

La \omterm{def:qm:bayesian-methods:predictive}{loi prédictive a posteriori} du \omterm{def:qm:estimation:sharpe}{ratio de Sharpe} des trois prochaines années du meilleur gérant est une loi normale de moyenne 1.14 et
d'écart type $\sqrt{0.33^2 + 1/3} = 0.67$~: une probabilité de 4.3~\% d'avoir une période négative, et une probabilité de 2.9~\% de retrouver le 2.41 qui l'a fait remarquer.

\section{Monte-Carlo par chaînes de Markov}

Le \omterm{def:qm:bayesian-methods:hier}{bayésien empirique} injecte $\hat\tau$ comme s'il était connu~; avec cinquante unités, $\tau$ est incertain, et un traitement bayésien complet l'intègre
sur lui. La \omterm{def:qm:bayesian-methods:prior}{loi a posteriori} n'est alors plus normale, et se calcule par simulation.

\begin{definition}[Monte-Carlo par chaînes de Markov, Metropolis--Hastings, échantillonneur de Gibbs]\label{def:qm:bayesian-methods:mcmc}
Le \emph{Monte-Carlo par chaînes de Markov}\index{Monte-Carlo par chaînes de Markov} (MCMC) tire des valeurs d'une densité $\pi$ connue à une constante près, en faisant tourner une
\omterm{def:qm:markov-chains-and-queues:ctmc}{chaîne de Markov} dont la \omterm{def:qm:stochastic-differential-equations:kolmogorov}{loi stationnaire} est $\pi$. L'\emph{algorithme de Metropolis--Hastings}\index{algorithme de Metropolis--Hastings}
propose $y \sim q(\cdot \mid x)$ à partir de l'état courant $x$ et l'accepte avec la probabilité $\alpha(x, y) = \min\bigl(1,
\frac{\pi(y)q(x \mid y)}{\pi(x)q(y \mid x)}\bigr)$, en restant en $x$ sinon~; avec une proposition symétrique de marche aléatoire, le rapport vaut
$\pi(y)/\pi(x)$. L'\emph{échantillonneur de Gibbs}\index{echantillonneur@échantillonneur de Gibbs} met à jour un bloc de coordonnées à la fois, par un tirage selon sa loi conditionnelle
sachant les autres.
\end{definition}

\begin{proposition}[Équilibre détaillé de Metropolis--Hastings]\label{prop:qm:bayesian-methods:mh}
Le noyau de Metropolis--Hastings vérifie $\pi(x)K(x, y) = \pi(y)K(y, x)$ pour $x \ne y$, donc $\pi$ est stationnaire. Une mise à jour de Gibbs est un
pas de Metropolis--Hastings dont la proposition est la loi conditionnelle, et sa probabilité d'acceptation vaut 1.
\end{proposition}

\begin{proof}
Pour $x \ne y$, $\pi(x)K(x, y) = \pi(x)q(y \mid x)\alpha(x, y) = \min\bigl(\pi(x)q(y \mid x), \pi(y)q(x \mid y)\bigr)$, qui est symétrique en
$(x, y)$. En sommant l'\omterm{def:qm:markov-chains-and-queues:balance}{équilibre détaillé} sur $x$, on obtient $\sum_x\pi(x)K(x, y) = \pi(y)$ (l'analogue à temps discret de la condition
du chapitre~8). Pour Gibbs, si $x$ et $y$ ne diffèrent que par le bloc $j$, $q(y \mid x) = \pi(y_j \mid x_{-j})$ et $x_{-j} = y_{-j}$, donc
$\pi(y)q(x \mid y) = \pi(x_{-j})\pi(y_j \mid x_{-j})\pi(x_j \mid x_{-j}) = \pi(x)q(y \mid x)$, et le rapport vaut 1.
\end{proof}

\begin{definition}[Période de chauffe (burn-in), taille d'échantillon effective]\label{def:qm:bayesian-methods:ess}
La \emph{période de chauffe}\index{période de chauffe} (burn-in) est la portion initiale d'une chaîne, écartée parce qu'elle dépend encore du point de départ. La
\emph{taille d'échantillon effective}\index{taille d'échantillon effective} (effective sample size) de $n$ tirages corrélés est $n/(1 + 2\sum_{k \ge 1}\rho_k)$, $\rho_k$ étant les
autocorrélations de la chaîne~: le nombre de tirages indépendants ayant la même variance de la moyenne.
\end{definition}

Pour le \omterm{def:qm:bayesian-methods:hier}{modèle hiérarchique} à \omterm{def:qm:bayesian-methods:prior}{lois a priori} plates sur $m$ et $\tau$, toutes les lois conditionnelles sont usuelles~: chaque $\theta_i$ est normal (d'après la
proposition précédente), $m$ est normal autour de $\bar\theta$, et $\tau^2$ suit une loi gamma inverse. Un \omterm{def:qm:bayesian-methods:mcmc}{échantillonneur de Gibbs} de 20\,000 balayages (2\,000 de \omterm{def:qm:bayesian-methods:ess}{période de chauffe}) donne au
meilleur gérant une moyenne a posteriori de 1.12, avec un \omterm{def:qm:bayesian-methods:predictive}{intervalle de crédibilité} à 95~\% $(0.41, 2.01)$, plus large que celui du \omterm{def:qm:bayesian-methods:hier}{bayésien empirique}, $1.14 \pm 1.96 \times 0.33$
parce que $\tau$ est incertain~: sa moyenne a posteriori vaut 0.40, avec un intervalle à 95~\% $(0.13, 0.66)$. Un échantillonneur de Metropolis par marche aléatoire, appliqué à la
\omterm{def:qm:bayesian-methods:prior}{loi a posteriori} marginale de $(m, \ln\tau)$, quatre chaînes issues de points de départ dispersés, accepte 50~\% de ses propositions et concorde (moyenne a posteriori de $\tau$
0.39, voir la \cref{fig:qm:bayesian-methods:tau})~; sa statistique de Gelman--Rubin vaut 1.001, et les 18\,000 tirages conservés de $\tau$ valent 506
tirages indépendants de la chaîne de Gibbs et 767 d'une seule chaîne de Metropolis.

\begin{omfigure}
\begin{tikzpicture}
\begin{axis}[width=10cm,height=5cm,xlabel={$\tau$, dispersion des vrais ratios de Sharpe},ylabel={densité a posteriori},
  tick label style={font=\small},label style={font=\small},xmin=0,xmax=1,ymin=0,ymax=3.6,grid=major,grid style={gray!25},
  legend columns=3,legend style={font=\footnotesize,at={(0.5,-0.3)},anchor=north,draw=none,fill=none,
  /tikz/every even column/.append style={column sep=8pt}},table/col sep=comma]
\addplot[omThm,thick,const plot mark mid,fill=omThm!15] table[x=tau,y=gibbs]{figdata/methods/14-bayesian-methods/tau.csv};
\addplot[omDef,very thick,const plot mark mid] table[x=tau,y=metropolis]{figdata/methods/14-bayesian-methods/tau.csv};
\addplot[black,dashed,thin,sharp plot] coordinates {(0.41,0) (0.41,3.6)};
\legend{échantillonneur de Gibbs,Metropolis (quatre chaînes),bayésien empirique $\hat\tau$}
\end{axis}
\end{tikzpicture}

\omcaption{\omterm{def:qm:bayesian-methods:prior}{Loi a posteriori} de la dispersion $\tau$ des vrais \omterm{def:qm:estimation:sharpe}{ratios de Sharpe} de la plateforme, issue de l'\omterm{def:qm:bayesian-methods:mcmc}{échantillonneur de Gibbs} (18\,000 tirages) et de quatre
chaînes de Metropolis à marche aléatoire sur la loi marginale a posteriori, face à l'estimation ponctuelle du \omterm{def:qm:bayesian-methods:hier}{bayésien empirique} 0.41 (en tirets). Données~: le tutoriel du chapitre,
avec graine fixée.}\label{fig:qm:bayesian-methods:tau}
\end{omfigure}

\section{Tutoriel~: la liste restreinte de l'allocateur}\label{tut:qm:bayesian-methods:shortlist}

\begin{tutorial}
\textbf{Objectif.} Faire rétrécir vers la moyenne les \omterm{def:qm:estimation:sharpe}{ratios de Sharpe} de cinquante gérants par la méthode bayésienne empirique et par un \omterm{def:qm:bayesian-methods:mcmc}{échantillonneur de Gibbs}, et vérifier ce rétrécissement sur les trois
années suivantes. \textbf{État final~:} les \cref{fig:qm:bayesian-methods:shrink,fig:qm:bayesian-methods:mse,fig:qm:bayesian-methods:tau}, le meilleur
gérant, ramené de 2.41 à 1.14, et la baisse de 32~\% de l'erreur hors échantillon.
\begin{tutsteps}
\item \textbf{\omterm{def:qm:bayesian-methods:hier}{Bayésien empirique}} pour des moyennes normales, avec des \omterm{def:qm:estimation:estimator}{erreurs types} égales ou inégales.
\omcode{code/firm/bayes/firm_bayes.py}{45}{66}{Rétrécissement bayésien empirique de moyennes normales.}\label{lst:qm:bayesian-methods:eb}
\item \textbf{L'\omterm{def:qm:bayesian-methods:mcmc}{échantillonneur de Gibbs}} pour le modèle normal hiérarchique.
\omcode{code/firm/bayes/firm_bayes.py}{96}{113}{Échantillonneur de Gibbs pour $x_i \sim \mathcal N(\theta_i, \sigma_i^2)$, $\theta_i \sim \mathcal N(m,
\tau^2)$.}\label{lst:qm:bayesian-methods:gibbs}
\item \textbf{Exécuter} \texttt{problem()}, \texttt{average\_mse()}, \texttt{mcmc()} dans \texttt{qm\_bayes.py}, puis \texttt{fig\_bayes.py}.
\end{tutsteps}
\textbf{Pour aller plus loin.} Donner aux gérants des longueurs d'historique différentes (des $\sigma_i$ inégaux) et observer comment les historiques courts
sont davantage rétrécis~; remplacer la loi normale a priori par une loi de Student $t$, de sorte qu'un gérant véritablement exceptionnel soit moins rétréci, et l'échantillonner avec la
routine de Metropolis.
\end{tutorial}

\section{Construction~: la boîte à outils bayésienne}\label{bld:qm:bayesian-methods:bayes}

\begin{build}
\bfield{Rôle} Chaque classement de nombreuses estimations bruitées dans la mini-firme (gérants, signaux, plateformes, contreparties) est rétréci avant
d'être utilisé~; ce module effectue le rétrécissement et l'échantillonnage.
\bfield{Interface} \texttt{beta\_binomial}, \texttt{normal\_normal}, \texttt{nig\_update}~; \texttt{eb\_normal\_means(x, se)} renvoyant le
$m$ et le $\tau^2$ a priori, les facteurs de rétrécissement et les moyennes et écarts types a posteriori~; \texttt{james\_stein(x, se)}~;
\texttt{metropolis(logpdf, x0, n, step, rng)}~; \texttt{gibbs\_hierarchical(x, se, n\_iter, rng)}~; \texttt{ess(chain)}~;
\texttt{rhat(chains)}.
\bfield{Règles} Les \omterm{def:qm:bayesian-methods:prior}{lois a priori} et leurs paramètres sont indiqués avec chaque \omterm{def:qm:bayesian-methods:prior}{loi a posteriori}~; les échantillonneurs prennent un générateur et renvoient la chaîne entière, et aucun
résumé a posteriori n'est donné sans une \omterm{def:qm:bayesian-methods:ess}{taille d'échantillon effective} et, pour plusieurs chaînes, $\hat R$.
\bfield{Tests d'acceptation} \texttt{code/firm/bayes/tests/}~: les formules conjuguées~; l'estimation bayésienne empirique retrouve $(m, \tau)$ et son point fixe
maximise la vraisemblance marginale avec des erreurs inégales~; James--Stein fait mieux que les estimations brutes~; Metropolis échantillonne une loi normale~; la \omterm{def:qm:bayesian-methods:ess}{taille
d'échantillon effective} d'une chaîne AR(1) vaut $n(1 - \phi)/(1 + \phi)$~; Gibbs concorde avec le \omterm{def:qm:bayesian-methods:hier}{bayésien empirique} lorsque les unités sont nombreuses.
\bfield{Pour aller plus loin} Une \omterm{def:qm:bayesian-methods:prior}{loi a priori} de Student $t$ par Metropolis dans Gibbs~; le Monte-Carlo hamiltonien sur la \omterm{def:qm:bayesian-methods:prior}{loi a posteriori} marginale~; la formule de Robbins pour
les comptages de Poisson.
\end{build}

\begin{omsources}
\item T.~Bayes, ``An essay towards solving a problem in the doctrine of chances'', \emph{Philosophical Transactions of the Royal Society}
53, 1763.
\item W.~James and C.~Stein, ``Estimation with quadratic loss'', \emph{Fourth Berkeley Symposium}, vol.~1, 1961.
\item B.~Efron and C.~Morris, ``Data analysis using Stein's estimator and its generalizations'', \emph{Journal of the American Statistical
Association} 70, 1975.
\item N.~Metropolis, A.~W.~Rosenbluth, M.~N.~Rosenbluth, A.~H.~Teller and E.~Teller, ``Equation of state calculations by fast computing
machines'', \emph{Journal of Chemical Physics} 21, 1953.
\item W.~K.~Hastings, ``Monte Carlo sampling methods using Markov chains and their applications'', \emph{Biometrika} 57, 1970.
\item S.~Geman and D.~Geman, ``Stochastic relaxation, Gibbs distributions, and the Bayesian restoration of images'', \emph{IEEE
Transactions on Pattern Analysis and Machine Intelligence} 6, 1984.
\item A.~Gelman and D.~B.~Rubin, ``Inference from iterative simulation using multiple sequences'', \emph{Statistical Science} 7, 1992.
\end{omsources}

\section{Exercices}

\begin{exercise}[$\star$]\label{exo:qm:bayesian-methods:1}
Un nouveau signal a été juste sur 30 transactions sur 50. Avec une \omterm{def:qm:bayesian-methods:prior}{loi a priori} $\mathrm{Beta}(50, 50)$, déterminer la \omterm{def:qm:bayesian-methods:prior}{loi a posteriori} et sa moyenne~; refaire le calcul avec une \omterm{def:qm:bayesian-methods:prior}{loi a priori} plate
$\mathrm{Beta}(1, 1)$.
\end{exercise}

\begin{exercise}[$\star$]\label{exo:qm:bayesian-methods:2}
Le \omterm{def:qm:estimation:sharpe}{ratio de Sharpe} sur cinq ans d'un gérant est de 1.5. Avec une \omterm{def:qm:bayesian-methods:prior}{loi a priori} $\mathcal N(0.5, 0.4^2)$ pour les vrais \omterm{def:qm:estimation:sharpe}{ratios de Sharpe} et une \omterm{def:qm:estimation:estimator}{erreur type} de
$1/\sqrt5$, déterminer la moyenne et l'écart type a posteriori.
\end{exercise}

\begin{exercise}[$\star$]\label{exo:qm:bayesian-methods:3}
Une chaîne de 10\,000 tirages a une autocorrélation d'ordre un égale à 0.8 et se comporte comme un AR(1). Calculer sa \omterm{def:qm:bayesian-methods:ess}{taille d'échantillon effective}.
\end{exercise}

\begin{exercise}[$\star\star$]\label{exo:qm:bayesian-methods:4}
Montrer que, avec $k$ unités, un $\sigma$ commun et l'\omterm{def:qm:estimation:estimator}{estimateur} des moments $\hat\tau^2 = s_x^2 - \sigma^2 > 0$, le facteur de rétrécissement \omterm{def:qm:bayesian-methods:hier}{bayésien empirique} est
$\sigma^2/s_x^2$, et le comparer au facteur de James--Stein vers la moyenne globale.
\end{exercise}

\begin{exercise}[$\star\star$]\label{exo:qm:bayesian-methods:5}
Dans l'exemple de la plateforme, déterminer combien d'années d'historique feraient tomber le facteur de rétrécissement $B$ à la moitié.
\end{exercise}

\begin{exercise}[$\star\star$]\label{exo:qm:bayesian-methods:6}
Établir la loi conditionnelle de $\tau^2$ dans l'\omterm{def:qm:bayesian-methods:mcmc}{échantillonneur de Gibbs} lorsque la \omterm{def:qm:bayesian-methods:prior}{loi a priori} sur $\tau$ est plate.
\end{exercise}

\begin{exercise}[$\star\star\star$]\label{exo:qm:bayesian-methods:7}
\emph{Programmation.} Sur 2\,000 plateformes simulées, comparer la moyenne du meilleur historique de trois ans, le vrai \omterm{def:qm:estimation:sharpe}{ratio de Sharpe} du gérant choisi, son
estimation bayésienne empirique et ses trois années suivantes.
\end{exercise}

\begin{exercise}[$\star\star\star$]\label{exo:qm:bayesian-methods:8}
\emph{Trouver l'erreur.} «~Nous rétrécissons vers zéro, par le facteur de James--Stein, le \omterm{def:qm:estimation:sharpe}{ratio de Sharpe} de backtest de chaque signal, si bien que nos estimations sont
prudentes et que nous n'avons plus à nous soucier du nombre de signaux testés.~»
\end{exercise}

\section{Problème~: la liste restreinte de l'allocateur}

\begin{problem}[{Problème du week-end --- rétrécissement du meilleur de cinquante gérants}]\label{pb:qm:bayesian-methods:1}
Une plateforme compte cinquante gérants dont les vrais \omterm{def:qm:estimation:sharpe}{ratios de Sharpe} annuels sont tirés de $\mathcal N(0.5, 0.4^2)$~; chacun dispose d'un historique de trois ans avec
une \omterm{def:qm:estimation:estimator}{erreur type} de $1/\sqrt3$ et, plus tard, un second historique de trois ans. La plateforme est celle du chapitre, graine fixée.

\medskip\noindent\textbf{Partie I --- Les historiques.}
\begin{enumerate}
\item Quel est le meilleur historique, et quelles sont la moyenne et l'écart type des cinquante historiques~?
\item Combien de gérants ont un historique négatif~?
\item Quel écart type les historiques auraient-ils si tous les gérants avaient le même vrai \omterm{def:qm:estimation:sharpe}{ratio de Sharpe}~?
\item Quelles sont les estimations des moments de $m$ et de $\tau$~?
\item Quel est le facteur de rétrécissement $B$~?
\end{enumerate}

\medskip\noindent\textbf{Partie II --- Rétrécissement.}
\begin{enumerate}[resume]
\item Quelles sont la moyenne et l'écart type a posteriori du meilleur gérant~?
\item Que donne James--Stein~?
\item Que deviennent les deuxième et troisième historiques, et le classement change-t-il~?
\item Quels étaient les vrais \omterm{def:qm:estimation:sharpe}{ratios de Sharpe} des trois premiers, et que valent leurs trois années suivantes~?
\item Quelle est la probabilité prédictive a posteriori que les trois années suivantes du meilleur gérant soient négatives~?
\end{enumerate}

\medskip\noindent\textbf{Partie III --- Vérification.}
\begin{enumerate}[resume]
\item De combien le \omterm{def:qm:bayesian-methods:hier}{bayésien empirique} réduit-il l'erreur par rapport aux vrais \omterm{def:qm:estimation:sharpe}{ratios de Sharpe} sur cette plateforme~?
\item Et en moyenne sur 2\,000 plateformes, par rapport à la vérité et aux trois années suivantes~?
\item En moyenne, que valent le meilleur historique, le \omterm{def:qm:estimation:sharpe}{ratio de Sharpe} réel du gérant choisi et son estimation rétrécie~?
\item Que disent la moyenne a posteriori et l'intervalle à 95~\% de l'\omterm{def:qm:bayesian-methods:mcmc}{échantillonneur de Gibbs} pour le meilleur gérant~?
\item Quelle est la \omterm{def:qm:bayesian-methods:prior}{loi a posteriori} de $\tau$, et pourquoi importe-t-elle~?
\end{enumerate}

\medskip\noindent\textbf{Partie IV --- Jugement.}
\begin{enumerate}[resume]
\item Le comité doit-il doubler le capital du meilleur gérant~?
\item Quelle \omterm{def:qm:bayesian-methods:prior}{loi a priori} choisir pour un gérant recruté hors de la plateforme~?
\item Dans quel cas la \omterm{def:qm:bayesian-methods:prior}{loi a priori} normale rétrécit-elle trop un gérant véritablement exceptionnel~?
\item Énoncer le \emph{résultat nommé}~: la moyenne a posteriori du meilleur gérant et la réduction hors échantillon de l'erreur.
\item En une phrase~: pourquoi faut-il rétrécir le meilleur de nombreux historiques bruités~?
\end{enumerate}
\end{problem}

\section{Questions d'entretien}

\begin{interviewq}[$\star$ \iqroles{chercheur, trader}]\label{iq:qm:bayesian-methods:1}
Quelle est la différence entre un intervalle de confiance et un \omterm{def:qm:bayesian-methods:predictive}{intervalle de crédibilité}~?
\end{interviewq}

\begin{interviewq}[$\star\star$ \iqroles{chercheur, trader}]\label{iq:qm:bayesian-methods:2}
Vous devez choisir l'un des cinquante gérants d'après leurs \omterm{def:qm:estimation:sharpe}{ratios de Sharpe} sur trois ans. Comment estimeriez-vous ce que le meilleur d'entre eux rapportera~?
\end{interviewq}

\begin{interviewq}[$\star\star$ \iqroles{chercheur, ingénieur ML}]\label{iq:qm:bayesian-methods:3}
Pourquoi l'\omterm{def:qm:bayesian-methods:js}{estimateur de James--Stein} peut-il battre les moyennes empiriques, et quel est le piège~?
\end{interviewq}

\begin{interviewq}[$\star\star$ \iqroles{chercheur, ingénieur ML}]\label{iq:qm:bayesian-methods:4}
Comment savez-vous qu'un échantillonneur MCMC a convergé~?
\end{interviewq}

\begin{interviewq}[$\star\star$ \iqroles{ingénieur ML}]\label{iq:qm:bayesian-methods:5}
Comment la régression ridge se rattache-t-elle à une \omterm{def:qm:bayesian-methods:prior}{loi a priori} bayésienne~?
\end{interviewq}

\begin{interviewq}[$\star\star\star$ \iqroles{chercheur}]\label{iq:qm:bayesian-methods:6}
Pourquoi l'\omterm{def:qm:bayesian-methods:mcmc}{algorithme de Metropolis--Hastings} fonctionne-t-il, et pourquoi un pas de Gibbs est-il toujours accepté~?
\end{interviewq}
